\begin{table}[!t]
    \centering
\small
    \setlength{\tabcolsep}{5pt}
    \begin{tabular}{@{}M{0.215\linewidth} M{0.735\linewidth}@{}}
        \toprule
        \textbf{Element} & \textbf{Reading for this system} \\
        \midrule
        Adversary capabilities & No model access of any kind and no knowledge of the training corpus, only the ability to drive a periodic load on the photovoltaic plant, which an inverter switching schedule or a switchable load already provides. \\
        \addlinespace[3pt]
        Attack surface & The raw telemetry stream before tokenisation, reached by acting on the physical process or on the data path. \\
        \addlinespace[3pt]
        Manipulation of telemetry & A small periodic component at a candidate frequency, which is legitimate plant behaviour and survives any plausibility check on the raw signal. \\
        \addlinespace[3pt]
        Operational consequence & The component is present in the plant and, if structural aliasing is confirmed, under-represented in the model, so dispatch and curtailment are decided on an incomplete picture, and a detector reading the forecast residual is looking through the same tokeniser that lost it. \\
        \addlinespace[3pt]
        Residual risk & The pre-registered mitigation raises overlap by shortening the stride at fixed patch size, so it thins the stride branch alone and $\mathcal{F}_{P}$ survives at every deployed configuration. Shrinking the patch instead thins $\mathcal{F}_{P}$, and with it $\mathcal{F}_{S}$ once $S \le P$ forces the stride down as well; but that raises both fundamentals rather than closing either comb, so the candidate set is relocated and thinned, not removed. \\
        \bottomrule
    \end{tabular}
\caption{The adversarial-masking threat model, read through the five elements of the NIST taxonomy of attacks on machine-learning systems\autocite{vassilevAdversarialMachineLearning2024}. The reading separates what the actor needs from what the actor changes, and both from the exposure that survives the mitigation this report pre-registers. Its distinguishing feature is the first row: the capability required is not access to the model but the ability to make the plant do something the plant is entitled to do, so nothing in the manipulation is anomalous at the point where a plausibility check would sit. \autoref{tab:riskRegister} carries this as its adversarial-masking row at high residual exposure.}
    \label{tab:threatModel}
\end{table}
